\section{Discussion}

\subsection{Answers to the research questions}

\paragraph{Q-A (dense optimisations).} With identical inputs, arithmetic, timing boundary and a single thread, reordering the loops to i-k-j makes the C baseline $1.5$--$1.7\times$ faster for $n\le500$, $2$--$2.3\times$ at $n=750$--$1000$ and $3.7$--$6.9\times$
for $n\ge1024$; blocking with $b=96$ adds nothing up to $n=1000$ and a further $1.4$--$1.7\times$ over i-k-j for $n\ge2000$, giving up to $9.9\times$ over the baseline at $n=2048$ (Table~\ref{tab:dense}, Figure~\ref{fig:speed}). The optimised kernels cost
$0.64$--$0.89$\,ns (i-k-j, $n\le1250$) or $0.68$--$0.85$\,ns (blocked, all sizes) per multiply--add and keep a near-cubic growth ($k=3.04\pm0.07$ for blocked), whereas the baseline's cost per multiply--add grows from $1.1$ to $8.0$\,ns. As a consequence the largest feasible size under the same budget moves from $3000$ to $6000$.
The speed-up numbers are bounded by the baseline's irregular behaviour (spikes at $512$ and $1024$): they should be read as ``at these sizes'' and not as a smooth law.

\paragraph{Q-B (sparse versus dense).} For inputs with $0.1$--$50\%$ nonzeros the own CSR$\times$CSR kernel is faster than the best dense method on the same matrices, by $10^5\times$ at $d=0.1\%$ (kernel only), $2\cdot10^2$--$4\cdot10^3\times$ at $1\%$ and $2$--$3\times$ at $50\%$; it is slower from a break-even density of about $72$--$87\%$ (precisely: $72$--$76\%$ at $n=2000$; $82$--$87\%$ at $n=500$ and $1000$, resolved to about $\pm10$ points).
The end-to-end advantage of the whole pipeline starting from dense matrices is limited by the dense$\to$CSR conversion ($O(n^2)$ and independent of $d$ to first order): at $n=1000$, $d=0.1\%$ it is $155\times$ instead of $10^5\times$. And the product of random sparse matrices is dense long before the inputs are: at $n=1000$, $d=5\%$ gives a $92\%$-dense $C$, and then
the cost per useful product ($9.7$\,ns at $5\%$, $30$\,ns at $1\%$) is $14$--$43\times$ that of a dense multiply--add ($0.7$\,ns; Table~\ref{tab:nsflop}). \emph{Where sparse helps}: very sparse inputs, banded or otherwise structured inputs whose product stays sparse, and inputs that are already stored in CSR. \emph{Where it stops helping}: dense outputs, inputs above $\approx70$--$80\%$ density, and
workloads that must convert from dense each time.

\paragraph{Separating the effects.} The optimisation effect (Part~A) was obtained with dense inputs, where the number of operations is fixed at $n^3$: changing the loop order or blocking changed the time by up to $9.9\times$ and nothing else. The storage and structure effect (Part~B) is of a different kind: it changes the \emph{number} of operations (from $n^3$ to $\approx d^2n^3$ useful products for random inputs),
by up to $10^5\times$, while the time per operation \emph{increases} ($1$ to $30$\,ns per product instead of $0.7$\,ns). The two are not additive: here the comparator in Part~B is the best optimised dense method, not the baseline; against the baseline all sparse speed-ups would be about $2\times$ ($n=1000$) to $6\times$ ($n=2000$) larger (baseline column of the tables). Also the
third factor, the \emph{compiler}, matters: with \texttt{-O3 -march=native} the dense blocked kernel became $3.3\times$ faster at $n=2000$, which would move the break-even density of Part~B down (the sparse kernel was not recompiled in the same comparison, so this is only an expectation).

\paragraph{Percentage of zeros is not enough.} At the same $d=1\%$ and $n=1000$, banded inputs run $15\times$ faster than random ones, produce a product $5\times$ sparser ($1.9\%$ versus $9.4\%$), and have $15\times$ lower cost per useful product. A statement such as ``a matrix with $99\%$ zeros is fast'' is therefore false in general: it depends on the pattern,
on the product's fill-in, on the format, on whether conversion is needed and on the implementation of the kernel.

\paragraph{Why a tile of 96 works although it does not fit in L1.} Three $96\times96$ tiles occupy $216$\,KiB, far more than the $32$\,KiB L1 data cache, yet $b=96$ was not slower than $b=32$ (which fits; $2.50$ versus $2.63$\,s, within the variability). The three-tiles-in-L1 criterion of Section~2 is conservative for this loop order: inside a tile the i-k-j order only needs one row of the $C$ tile and one row of the $B$ tile at a time ($768$\,bytes each) in L1, while the $B$ tile ($72$\,KiB) is reused from L2 for every row $i$ and is read with unit stride, a pattern the hardware prefetchers predict well. This explanation is a hypothesis; without hardware counters it was not measured.

\subsection{Hypotheses}
\begin{table}[H]
\centering\small
\begin{tabular}{@{}p{0.6cm}p{3.1cm}p{10.2cm}@{}}
\toprule
 & Verdict & Evidence \\
\midrule
H1 & supported & i-k-j is $1.45$--$6.9\times$ faster than i-j-k at every size; gain grows with $n$ but not monotonically; no spikes at $512$/$1024$. \\
H2 & partly supported & No gain of blocking for $n\le1000$ (as expected); a gain of $1.38$--$1.66\times$ for $n\ge2000$. The predicted bound ``$\lesssim1.5\times$'' is exceeded at $n=2000$ ($1.66\times$) and respected elsewhere. \\
H3 & supported & identical SHA-256 of $C$ in all $69$ groups; all tile sizes tested bit-for-bit equal. \\
H4 & supported, with a precise break-even & $10^5\times$ at $d=0.1\%$ down to $0.6$--$0.9\times$ at $90\%$; break-even $72$--$87\%$; random products fill in as $1-e^{-nd^2}$, banded ones stay banded. \\
H5 & supported & conversion is $92$--$99.9\%$ of the end-to-end time for banded inputs with $d\le1\%$; random $d=1\%$: $53$--$82\%$. \\
H6 & supported & dense: $24.00$ bytes per $n^2$; CSR: model within $1\%$ of the measured RSS at the largest scaling sizes. \\
\bottomrule
\end{tabular}
\caption{Predictions declared in the introduction against the results.}
\end{table}

\subsection{Strength of the evidence}
\begin{itemize}
  \item \textbf{Strong:} the ordering baseline $\ll$ i-k-j $\lesssim$ blocked for $n\ge2000$ (ratios $>3\times$ against samples whose IQR is at most $14\%$ of the median); the correctness of all results; the memory formulas; the orders of magnitude of the sparse advantage at low density, and the fact that it is reversed at $90\%$ for $n=2000$.
  \item \textbf{Moderate:} the break-even densities for $n\le1000$ (the speed-up is within $\pm15\%$ of $1$ over the last two densities); exponents (few points); the choice $b=96$ (flat landscape).
  \item \textbf{Weak or hypothesis:} all explanations in terms of caches, TLB, SIMD units or \texttt{realloc}; they are consistent with the data (e.g.\ the vectorisation evidence and the compiler experiment for SIMD) but no hardware counters were available. For the baseline's spikes at powers of two a cache-associativity conflict is the usual explanation \cite{drepper,hennessy}; this study neither tests nor excludes it.
\end{itemize}

\subsection{Limitations}
\begin{itemize}
  \item \textbf{One machine.} A virtual machine with $2$ vCPUs, $8$\,GiB, shared host, unpinned, no visible frequency governor. Between-process variability reaches $13$--$35\%$ at $n=1000$; with $2$--$5$ samples at large $n$ the medians are less certain than the IQR suggests. Results are \emph{not} from the author's laptop and would differ in absolute terms there.
  \item \textbf{No CPU pinning.} Pinning each process to one logical CPU (\texttt{taskset}) could reduce the between-process noise and will be used in Assignment~3 (threads); in a virtual machine it fixes the process to a vCPU, not to a physical core, so it can reduce but not remove the variability.
  \item \textbf{Few samples at the largest sizes} (budget rule): baseline $n=3000$ has $2$ samples; $4000$--$6000$ have $2$--$3$; the sparse scaling has $5$.
  \item \textbf{Implementation scope.} One sparse format (CSR) and one algorithm (Gustavson with a dense accumulator); no comparison with a library (SciPy was used only to validate), no CSC/COO, no tuned sort. A faster sparse kernel would move the break-even; a library dense kernel (BLAS) would very likely be faster than all three dense methods and would move it down.
  \item \textbf{Inputs.} Two structures only (random and banded) and positive values; real matrices (power-law degree, block structure) may behave differently. Matched comparisons only up to $n=2000$; the baseline is run on only two cells at $n=2000$.
  \item \textbf{Exploratory decisions.} A pilot at $n=1500$ and a smoke test at $n=500$ were seen before the hypotheses were fixed; the densities $70\%$ and $90\%$, the compiler experiment and the variability experiment were added after those looks. The tile was chosen on the same machine at a size outside the comparison.
  \item \textbf{Scaling runs} use CSR files, so they exclude conversion by design, and only random patterns.
\end{itemize}

\section{Conclusions}
\begin{enumerate}
  \item On dense $n\times n$ \texttt{float64} matrices in sequential C at \texttt{-O2}, the i-j-k baseline loses to i-k-j by $1.5\times$ to $6.9\times$, and to the blocked method ($b=96$) by up to $9.9\times$; blocking only pays beyond $n\approx1250$, where it adds $1.4$--$1.7\times$ over i-k-j. The largest feasible size within $300$\,s doubles, from $3000$ to $6000$. Results are bit-identical.
  \item Cost per multiply--add: $0.64$--$0.89$\,ns (i-k-j) and $0.68$--$0.85$\,ns (blocked) against $1.1$--$8.0$\,ns (baseline); exponent $3.04\pm0.07$ (blocked).
  \item The own CSR$\times$CSR kernel is $10^2$--$10^5\times$ faster than the best dense method for $d\le1\%$, still $2$--$3\times$ at $50\%$, and slower above a break-even density of about $70$--$85\%$ on this machine; a dense$\to$CSR conversion removes most of the advantage at high sparsity, and random products become dense after a few percent of input density.
  \item The percentage of zeros alone does not predict performance: at $1\%$, banded inputs are $15\times$ faster than random ones.
  \item With CSR inputs, $n=64000$ ($d=0.1\%$) and $n=16000$ ($d=1\%$) are feasible in $14$\,s and $4.3$\,s within $3$\,GiB; the next sizes were excluded by predicted memory ($>12$\,GiB), and no run failed. Dense matrices of $n=64000$ would need $98$\,GB.
  \item A compiler flag (vectorisation) had an effect comparable to or larger than the algorithmic changes in the dense part ($3.3\times$ for the blocked kernel at $n=2000$); this should be kept in mind when the next assignment adds vectorisation and parallelism.
\end{enumerate}

\section{Use of generative AI}
In line with the course policy, the use of AI is disclosed: the tool was \textbf{Claude (Anthropic)}. It was used for planning the experimental design, writing the C implementations, the input generator, the benchmark driver, the test, analysis and plotting scripts, the supplementary experiments and a first draft of this report.
The affected parts are \texttt{src/}, \texttt{scripts/} and \texttt{report/}. The author is responsible for understanding the code, validating the measurements, checking the references and defending the conclusions; generated text is not evidence, and every number in this report was produced by the included code
from the raw data in \texttt{results/}.

\section{Reproducibility materials}
\label{sec:materials}
Source code, build and run instructions, raw measurements, input-generation scripts and analysis scripts are in the repository \repourl{} and in \path{Assignment02_GonzalezFernandez_Materials.zip}. The \texttt{README.md} gives the exact commands
(\texttt{make build inputs test bench analyze report}); the input matrices are not stored (about $5$\,GB) but regenerated with \texttt{scripts/gen\_inputs.py} and verified against the SHA-256 digests in \texttt{data/manifest.json}; the dense inputs are identical to those of Assignment~1 (\texttt{data/a1\_reference\_sha256.json}).
Raw data: \path{results/raw_measurements.csv} (one row per sample, warm-up included) and \path{results/runs.csv} (one row per process); environment: \path{results/environment.json}; every figure and table is generated by \path{scripts/analyze.py}. The measurement driver reads \texttt{/proc} and therefore needs Linux;
the full campaign took several hours on the reported machine. The baseline kernel is that of Assignment~1 \cite{a1}, whose structure was taken from the example implementations on the course website \cite{course}; all other code, the driver, validation, supplementary experiments and analysis were written for this study with the AI assistance disclosed above.
No other external code, text or figures are reused.

\begin{thebibliography}{99}
\bibitem{a1} H.~Gonz\'alez Fern\'andez, ``Assignment 1: How runtime and memory scale in the basic triple-loop matrix multiplication in C, Java and Python,'' Big Data, Univ. Las Palmas de Gran Canaria, Oct. 2026. Code: \url{https://github.com/hugogzfz/bd-assignment1-matmul}, commit 101af1c.
\bibitem{course} ``Ejemplos'' (matrix-multiplication examples in Python, Java, Rust and C), \emph{Big Data} course website, Block~1, Complexity Management. \url{https://mltechdrawer.github.io/BIGDATA/Block1_Theoretical_Concepts_of_Big_Data/Complexity_Management/codigos/}. Accessed: Oct.~14, 2026.
\bibitem{lam} M.~D.~Lam, E.~E.~Rothberg and M.~E.~Wolf, ``The cache performance and optimizations of blocked algorithms,'' in \emph{Proc. 4th Int. Conf. Architectural Support for Programming Languages and Operating Systems (ASPLOS IV)}, Santa Clara, CA, USA, 1991, pp.~63--74.
\bibitem{goto} K.~Goto and R.~A.~van de Geijn, ``Anatomy of high-performance matrix multiplication,'' \emph{ACM Trans. Math. Softw.}, vol.~34, no.~3, pp.~1--25, 2008.
\bibitem{gustavson} F.~G.~Gustavson, ``Two fast algorithms for sparse matrices: Multiplication and permuted transposition,'' \emph{ACM Trans. Math. Softw.}, vol.~4, no.~3, pp.~250--269, Sep. 1978.
\bibitem{saad} Y.~Saad, \emph{Iterative Methods for Sparse Linear Systems}, 2nd~ed. Philadelphia, PA, USA: SIAM, 2003.
\bibitem{drepper} U.~Drepper, ``What every programmer should know about memory,'' Red~Hat, Inc., Nov. 2007.
\bibitem{hennessy} J.~L.~Hennessy and D.~A.~Patterson, \emph{Computer Architecture: A Quantitative Approach}, 6th~ed. Cambridge, MA, USA: Morgan Kaufmann, 2017.
\bibitem{higham} N.~J.~Higham, \emph{Accuracy and Stability of Numerical Algorithms}, 2nd~ed. Philadelphia, PA, USA: SIAM, 2002.
\bibitem{goldberg} D.~Goldberg, ``What every computer scientist should know about floating-point arithmetic,'' \emph{ACM Comput. Surv.}, vol.~23, no.~1, pp.~5--48, Mar. 1991.
\bibitem{georges} A.~Georges, D.~Buytaert and L.~Eeckhout, ``Statistically rigorous Java performance evaluation,'' in \emph{Proc. 22nd ACM SIGPLAN Conf. Object-Oriented Programming Systems, Languages and Applications (OOPSLA)}, Montreal, Canada, 2007, pp.~57--76.
\bibitem{hoefler} T.~Hoefler and R.~Belli, ``Scientific benchmarking of parallel computing systems: Twelve ways to tell the masses when reporting performance results,'' in \emph{Proc. Int. Conf. High Performance Computing, Networking, Storage and Analysis (SC '15)}, Austin, TX, USA, 2015, Art.~no.~73.
\bibitem{numpy} C.~R.~Harris \emph{et al.}, ``Array programming with NumPy,'' \emph{Nature}, vol.~585, pp.~357--362, 2020, doi: 10.1038/s41586-020-2649-2.
\bibitem{scipy} P.~Virtanen \emph{et al.}, ``SciPy 1.0: Fundamental algorithms for scientific computing in Python,'' \emph{Nature Methods}, vol.~17, pp.~261--272, 2020, doi: 10.1038/s41592-019-0686-2.
\end{thebibliography}
